\chapter{转换流程与实现方法}

整体流程划分为输入解析、语义归一、双格式生成和输出校验四层。
输入层定位主文件、元数据、章节、文献和图片；语义层提取标题、段落、图表、公式与引用关系；
生成层分别调用 LaTeX 和 DOCX 处理链；校验层检查文件存在性、页面结构、样式属性与未解析引用。

\section{外部入口与路径隔离}

构建器以主文件所在目录作为论文源根目录，并将模板类复制到输出目录下的隔离运行区。
外部入口缺少元数据、章节或图片时应立即失败，不得回退读取模板仓库内的同名占位文件。

\section{语义转换与完整性检查}

\subsection{特征表示}

设需要检查的结构项共有 $n$ 个，第 $i$ 项通过检查时 $I_i=1$，否则 $I_i=0$，则结构完整率写为

\begin{equation}
  C=\frac{1}{n}\sum_{i=1}^{n} I_i.
  \label{eq:completeness}
\end{equation}

\subsubsection{数据不足时的回退策略}

自动转换无法可靠解释未登记的复杂宏时，流程应保留 LaTeX 作为版式主输出，并将 DOCX 标记为需要人工复核，而不是静默删除内容。

\begin{figure}[htbp]
  \centering
  \includegraphics[width=0.92\textwidth]{document-pipeline.png}
  \caption{学术文档双格式构建与校验流程}
  \label{fig:document-pipeline}
\end{figure}

正文交叉引用测试：转换流程如图\ref{fig:document-pipeline}所示。
